\section{Model Matching} \label{sec:model_matching}

In the process of autonomous wind turbine inspection, being aware of the relative position of the wind turbine with regards to the UAVs own position is of utmost importance. Obtaining and tracking this information can be done in various ways, for example trilateration with ultrasound or ultra-wide bandwidth technologies, GPS or searching the point cloud with LiDAR. Specifically, the system in this project works on the basis of a fusion of GPS and IMU signals for global localization, and point cloud model matching for localization relative to the object being inspected, e.g. a wind turbine.\\
Model matching refers to the process of locating a specific 3D model within the point cloud captured by a sensor system, e.g. a LiDAR sensor. It involves comparing the acquired point cloud data with a reference 3D model to identify the presence and precise location of the model within the scanned area. By finding correspondences between two point clouds, one can determine the transformation (position and orientation) required to align the model with the point cloud data.\\
The rest of the section presents the techniques used to perform model matching for this setting.


\subsection{3D Model generation}
Accurate and detailed 3D models are crucial for performing inspection of rotating wind turbines. However, while for the inspection of stationary wind turbines obtaining an exact model is relatively simple, finding a rotating wind turbine in the received point clouds presents a different problem. Many approaches had been used during the development of this system, but the best results were obtained by approximating the rotor in motion with a thin cylinder/plate model. This is a geometrical representation of the trace of the wind turbine's rotor during the rotation picked up by the LiDAR sensor.\\
For the purposes of the inspection workflow, it was desirable to have the system as modular and generic as possible, so the process of model generation was scripted in a Blender Python script file. This allows parametrization with regards to the wind turbine height and rotor radius.\\
% about blender and blender python scripting
Blender is a free and open source software tool developed for 3D computer graphics. The platform is used in the creation of animated films, visual effects, artwork, 3D printable models, motion graphics, etc. by a large community. In line with its open source nature, Blender offers an open library for Python scripting, making modelling through code possible.\\
The necessary input into the script is the wind turbine pillar height, the rotor radius, and its width. The script generates a 3d model of the rotating wind turbine (a cylinder approximating the rotor) in the Blender UI. The model can easily be exported to .obj and transformed in a point cloud ready for usage.


\subsection{Matching algorithm}
Figures~\ref{fig:matching_1},~\ref{fig:matching_2},~\ref{fig:matching_3}, and~\ref{fig:matching_4} illustrate the matching process as it appears in the RVIZ software. Figure~\ref{fig:matching_1} presents the initialized system with the wind turbine 3D model randomly positioned depicted in yellow and then fixed to the first "match" (depicted with red points) which is a group of points collected on a bush nearby in~\ref{fig:matching_2}. Figure~\ref{fig:matching_3} presents the next important situations in the matching process. Depicting the accumulated point cloud which now contains also points which are clearly from the actual wind turbine. Lastly, in figure~\ref{fig:matching_4}, once enough points coming from the wind turbine are accumulated, the matching process correctly matches the entire model (the yellow and red point clouds coincide with one another).

\begin{figure}
    \centering
    \includegraphics[width=0.45\textwidth]{figures/inspection_simulation/inspection_new/matching_initialization.png}
    \caption{Matching process - Initialization}
    \label{fig:matching_1}
\end{figure}
\begin{figure}
    \centering
    \includegraphics[width=0.45\textwidth]{figures/inspection_simulation/inspection_new/matching_first.png}
    \caption{Matching process - First match}
    \label{fig:matching_2}
\end{figure}
\begin{figure}
    \centering
    \includegraphics[width=0.45\textwidth]{figures/inspection_simulation/inspection_new/matching_accumulation.png}
    \caption{Matching process - Point cloud accumulation}
    \label{fig:matching_3}
\end{figure}
\begin{figure}
    \centering
    \includegraphics[width=0.45\textwidth]{figures/inspection_simulation/inspection_new/matching_full.png}
    \caption{Matching process - Full match}
    \label{fig:matching_4}
\end{figure}

\begin{algorithm}
    \caption{Matching algorithm pseudocode}\label{alg:matching}
    \begin{algorithmic}[1]
        \Procedure{Rotating wind turbine matching}{}
        \State Initialize drone takeoff
        \While{unsatisfactory match}
            \State Collect point cloud
            \State Point cloud registration
            \State Point cloud voxelization
            \State Ground plane rejection
            \State Point cloud clustering
            \For{i in clusters}
                \If{found better match}
                    \State Update best match
                \EndIf
            \EndFor
        \EndWhile\\
        
        \State \Return result
        \EndProcedure
    \end{algorithmic}
\end{algorithm}


The pseudocode for the model matching process used to align and match a 3D point cloud model of a rotating wind turbine during UAV drone inspection is given with~\ref{alg:matching}. The important steps of the algorithm are the following:
\begin{itemize}
    \item \textbf{Initialization} The process is initialized by placing a UAV drone in front of the wind turbine, clearing the site, and starting the automated inspection process. The UAV drone then takes off and starts collecting point clouds as it climbs in height.
    \item \textbf{Point Cloud Registration} Every new point cloud is incorporated into the global point cloud through the process called registration.
    \item \textbf{Point Cloud Voxelization}
    \item \textbf{Ground Plane Rejection} The points in the point cloud belonging to the ground plane are removed through the ground plane rejection algorithm. The algorithm searches for a plane with most points in its neighborhood through an optimization process.
    \item \textbf{Point cloud Clustering} The remaining point cloud is then segmented into clusters based on proximity. Clustering groups together points that are spatially close to each other, representing different objects in the scene.
    \item \textbf{Match Evaluation} Iterating over the clusters, the algorithm evaluates the match with the generated 3d model of the rotating wind turbine. The evaluation is done with the Iterative Closest Point (ICP) algorithm, keeping the pitch and roll values constant. This follows the assumption that the UAV drone stays upright during inspection.
\end{itemize}

Once the UAV drone reaches the desired height at which the inspection process can begin, the best match the algorithm found should very closely represent the actual position and orientation of the wind turbine relative to the drone.
